% =============================================================================
% Appendix: knowledge-area annotation --- design, distribution and consistency
% audit.
%
% Usage: \input this file after \appendix in the paper.
% Required packages: graphicx, booktabs, amsmath, amssymb.
% Assumes UTF-8 input and T1 font encoding (both provided by the ACL template).
%
% The annotation step was run once, with one model, one prompt and one
% taxonomy; it was not ablated and no accuracy was measured. Every number
% below comes from the artifacts of that run in output/knowledge-classification/
% (manifest.json, the run checkpoint and the published dataset) and from
% knowledge_classification.log. Figures are generated by
% ablations/knowledge_classification/make_figures.py, which also prints every
% number quoted below, checks the checkpoint against the published dataset and
% re-derives every macro-area from src/mmlu_pt/classification/taxonomy.json.
% By default figures are read from ./figures; to change this, define the macro
% before the \input, e.g.
%   \newcommand{\KnowledgeFigDir}{ablations/knowledge_classification/figures}
% =============================================================================
\providecommand{\KnowledgeFigDir}{figures}

\section{Knowledge-Area Annotation: Design, Distribution and Consistency Audit}
\label{app:kc}

Every released question carries two labels: a \emph{subject}, chosen from a
closed list specific to the question's academic level, and a
\emph{macro-area}, derived from the subject. The paper uses them to report
results per area and to stratify the sampling of few-shot demonstrations by
macro-area and level. The labels are assigned after semantic deduplication
(Appendix~\ref{app:dedup}) by a single pass of an open-weight language model
over the 37,887 retained questions, under constrained decoding that only
admits the subjects of the question's level (Table~\ref{tab:kc-overview}).
Unlike the filtering and deduplication stages, this step was run once, with
one model, one prompt and one taxonomy version: it was not ablated, and no
accuracy or human agreement was measured. This appendix therefore has a
different character from the preceding ones. It fixes the notation
(\S\ref{app:kc-setup}), records the design decisions and their motivation
(\S\ref{app:kc-taxonomy}, \S\ref{app:kc-prompt}), describes the labels that
resulted (\S\ref{app:kc-distribution}), audits their consistency against
metadata that the model never saw (\S\ref{app:kc-audit}), shows what the
labels look like at the boundaries between neighbouring subjects
(\S\ref{app:kc-boundary}) and reports cost and reproducibility
(\S\ref{app:kc-cost}). Limitations are discussed in
\S\ref{app:kc-limitations}.

\begin{table}[t]
\centering
\footnotesize
\setlength{\tabcolsep}{4pt}
\begin{tabular}{@{}lrrr@{}}
\toprule
Macro-area & Subjects & Questions & Share \\
\midrule
\multicolumn{4}{@{}l}{\emph{High school (8,746 questions, 17 subjects)}} \\
STEM & 7 & 4,964 & 56.76 \\
Languages and Arts & 5 & 2,279 & 26.06 \\
Humanities and Social Sciences & 4 & 1,479 & 16.91 \\
Health Sciences & 1 & 24 & 0.27 \\
\midrule
\multicolumn{4}{@{}l}{\emph{Undergraduate (29,141 questions, 55 subjects)}} \\
Law & 1 & 11,013 & 37.79 \\
STEM & 12 & 5,149 & 17.67 \\
Business and Economics & 5 & 4,116 & 14.12 \\
Health Sciences & 11 & 3,989 & 13.69 \\
Humanities and Social Sciences & 11 & 2,331 & 8.00 \\
Languages and Arts & 7 & 1,442 & 4.95 \\
Education & 1 & 433 & 1.49 \\
Agricultural and Veterinary Sciences & 3 & 417 & 1.43 \\
Services & 4 & 251 & 0.86 \\
\bottomrule
\end{tabular}
\caption{Macro-areas of taxonomy v1.0.0 and the questions assigned to each.
\emph{Subjects}: subjects of the macro-area at that level; \emph{Share}:
percentage of the level's questions. All 72 (level, subject) pairs received
at least one question. Over the whole set, Law holds 29.07\% of the
questions, STEM 26.69\%, Business and Economics 10.86\%, Health Sciences
10.59\%, Humanities and Social Sciences 10.06\% and Languages and Arts
9.82\%; Education, Agricultural and Veterinary Sciences and Services together
hold 2.90\%.}
\label{tab:kc-overview}
\end{table}

\paragraph{Summary.} Six decisions define the step. (1)~The model chooses
one subject from a closed, level-specific list, and the macro-area is a
function of the subject computed in code, so the two labels are always
consistent and the model never names a macro-area. (2)~The only evidence
shown to the model is the question and its choices: the answer key, the
exam, the edition, the course or job title and the source URL are withheld,
so that a label cannot be a shortcut from the exam's name. (3)~Decoding is
greedy, with the model's reasoning mode disabled, and constrained to the
allowed subject strings, so every record receives exactly one valid label and
no fallback, retry or ``Other'' category exists. (4)~The taxonomy is specific
to the project and written down with a definition, inclusions and exclusions
per subject and thirteen boundary rules; Law is a macro-area of its own,
Logic is separate from Mathematics and Computer Science, and Accounting
belongs to Business and Economics. (5)~Every output is validated against the
taxonomy before it is stored, results are checkpointed per batch with a hash
of each input row, and a run can only be resumed under an identical input,
model revision, prompt, taxonomy and code. (6)~The run took 84.9~minutes on
one GPU and processed 145.6~million prompt tokens. None of these choices was
compared with an alternative; the audit of \S\ref{app:kc-audit} measures the
consistency of the labels with withheld metadata, not their accuracy.

% -----------------------------------------------------------------------------
\subsection{Data, notation and metrics}
\label{app:kc-setup}

\paragraph{Data.} Let $D_S$ be the released set after semantic
deduplication: 37,887 records from 17 exams and 1,039 editions, 8,746 at
high-school level and 29,141 at undergraduate level, identified by the
SHA-256 of its records, \texttt{322a7279}\ldots, which the run manifest
stores. A record $x=(q,C,a,e,\ell,m)$ has question $q$, ordered choices
$C=(c_1,\dots,c_k)$, answer index $a$, exam $e$, level
$\ell\in\{\mathrm{HS},\mathrm{UG}\}$ and metadata $m$ (edition, question
number and source URL). Only $q$, $C$ and $\ell$ are shown to the model; $a$,
$e$ and $m$ are not.

\paragraph{Taxonomy and labels.} For each level $\ell$ the taxonomy fixes a
list of subjects $T_\ell$ ($|T_{\mathrm{HS}}|=17$, $|T_{\mathrm{UG}}|=55$)
and a map $\mu_\ell\colon T_\ell\to A_\ell$ to macro-areas
($|A_{\mathrm{HS}}|=4$, $|A_{\mathrm{UG}}|=9$); the 72 pairs $(\ell,s)$ with
$s\in T_\ell$ are the only labels a record can receive. The annotation is a
function $s(x)=f(q,C,\ell)\in T_\ell$, and the pipeline stores the pair
$(s(x),\mu_\ell(s(x)))$. Operationally, $f$ is the greedy token path of the
language model under a grammar that accepts exactly the strings of $T_\ell$
(vLLM's \emph{choice} structured output); because the constraint is applied
token by token, this path need not coincide with the subject of highest
total probability, a nuance we return to in \S\ref{app:kc-limitations}.

\paragraph{Metrics.} The descriptive statistics are counts $n_{\ell,s}$ and
shares per level, subject, macro-area and exam, and the number of distinct
subjects used by each exam. The consistency audit considers groups of
records $G$ (an exam, or the questions of one ENADE course) for which an
\emph{expected set} $E_G\subseteq T_\ell$ was fixed by rule before the labels
were inspected, and reports
\begin{equation*}
c_G=\frac{|\{x\in G: s(x)\in E_G\}|}{|G|},
\end{equation*}
the share of labels inside the expected set, with a 95\% Wilson interval,
together with the most frequent subject outside $E_G$. $c_G$ is a measure of
consistency, not of accuracy: $E_G$ describes the group, not each question,
and a question of a law course may legitimately require Philosophy or
Economics under the criterion of \S\ref{app:kc-prompt}. Prompt and output
lengths are the token counts recorded in the run checkpoint for every record.

% -----------------------------------------------------------------------------
\subsection{Taxonomy design}
\label{app:kc-taxonomy}

\paragraph{Motivation.} The labels must support MMLU-style reporting by
subject while matching what Brazilian exams actually assess: a school
curriculum in the vestibular and olympiad exams, and a landscape of
professions and degree programmes (law, health professions, engineering,
agricultural and service degrees) in the professional and civil-service
exams. Existing classifications, such as the areas of the national
high-school exam or the official classification of fields of study, cover
one of the two worlds but not both, and none has a place for the
rule-following puzzles of the informatics olympiad.

\paragraph{Design.} The taxonomy therefore has two level-specific lists under
nine macro-areas (Table~\ref{tab:kc-subjects}). The lists share the school
subjects and differ in the professional fields that only appear in
undergraduate exams; keeping them separate prevents a high-school question
from being labelled with a profession and an undergraduate question from
being forced into a school subject. Four choices depart from the usual
groupings. \emph{Law is a macro-area of its own:} the bar exam (OAB), the
national magistracy exam (ENAM) and the legal content of civil-service exams
make it the largest field of the corpus (29\% of all questions), and nesting
it under the humanities would leave that macro-area without a usable meaning
for reporting or stratification. \emph{Logic is a subject separate from
Mathematics and Computer Science:} the Brazilian Informatics Olympiad (OBI)
consists mostly of puzzles that are solved from rules given in the statement
and require neither mathematical content nor computational knowledge.
\emph{Accounting belongs to Business and Economics}, next to Finance,
Economics and the two administration subjects, because the accountants'
exam (CFCES) and the civil-service exams mix these fields. \emph{Physical
Education sits under Health Sciences} at both levels. Every subject has a
one-sentence definition, a list of inclusions and a list of exclusions that
name the neighbouring subject to which an excluded topic belongs (for
Mathematics, probability and statistical interpretation are excluded in
favour of Statistics). Thirteen boundary rules, written in Portuguese and
included verbatim in the prompt, settle the conflicts that recur in these
exams:
\begin{enumerate}\setlength{\itemsep}{0pt}
\item Portuguese covers text interpretation and the workings of the
language; Literature requires literary knowledge. The language of the text
does not define the subject.
\item Mathematics covers mathematical concepts and operations; Statistics
probability and inference; Logic deduction from rules entirely given in the
statement; Computer Science requires external computational knowledge. A
self-contained OBI ordering or constraint problem is Logic, even when framed
with programs.
\item Education requires didactics, curriculum or learning. The content of a
discipline in a teachers' exam belongs to the discipline taught.
\item Law requires legal interpretation; Accounting covers records,
statements, forensic accounting and auditing; Finance investments and
financial decisions. Legislation cited only as context does not make a
question Law.
\item Economics covers economic phenomena and models; Public Administration
state management and policy implementation.
\item Geography covers spatial organization and society--nature relations;
Earth Sciences requires specialized geological, geophysical or atmospheric
mechanisms. At high school, physical geography remains Geography.
\item Biology covers general biological processes; Biomedical Sciences
analyses and biomedical techniques; Medicine diagnosis and clinical conduct
for an individual. All medical specialties remain Medicine.
\item Public Health covers populations, epidemiology, surveillance and health
systems; care for one patient is Medicine.
\item Engineering covers the design and operation of physical systems;
Computer Science the workings of computers and software, including software
engineering.
\item Food Science and Technology covers food production and preservation;
Gastronomy culinary technique; Nutrition nutritional needs, dietary
assessment and diet therapy.
\item Occupational Health and Safety covers the prevention and control of
occupational risks; clinical diagnosis remains Medicine and legal
interpretation remains Law.
\item Marketing, human resources, logistics and secretarial work belong to
Business Administration; plant production and forestry to Agronomy; every
branch of law to Law.
\item An interdisciplinary question receives the single subject whose
knowledge is central to the command and the choices, never a fixed label of
the exam.
\end{enumerate}

\paragraph{Decision.} Taxonomy v1.0.0 was frozen before the run. Its SHA-256
(\texttt{3b2cdd5e}\ldots) is recorded in the run manifest, and the pipeline
refuses to resume a checkpoint produced under a different taxonomy or
prompt. Two invariants are checked when the taxonomy is loaded: Law is absent
from the high-school list, and Law is the only subject of its macro-area.

\begin{table*}[t]
\centering
\footnotesize
\setlength{\tabcolsep}{4pt}
\begin{tabular}{@{}lrrr@{}}
\toprule
Subject & HS & UG & Total \\
\midrule
\multicolumn{4}{@{}l}{\emph{STEM}} \\
Mathematics & 1,489 & 664 & 2,153 \\
Statistics & 193 & 325 & 518 \\
Logic & 1,219 & 104 & 1,323 \\
Physics & 922 & 220 & 1,142 \\
Chemistry & 520 & 179 & 699 \\
Biology & 616 & 252 & 868 \\
Computer Science & 5 & 2,166 & 2,171 \\
Earth Sciences & -- & 71 & 71 \\
Environmental Sciences & -- & 220 & 220 \\
Engineering & -- & 558 & 558 \\
Architecture and Urbanism & -- & 117 & 117 \\
Food Science and Technology & -- & 273 & 273 \\
\midrule
\multicolumn{4}{@{}l}{\emph{Languages and Arts}} \\
Portuguese & 995 & 512 & 1,507 \\
English & 689 & 289 & 978 \\
Spanish & 5 & 172 & 177 \\
Literature & 488 & 104 & 592 \\
Arts & 102 & 126 & 228 \\
Linguistics & -- & 46 & 46 \\
Design & -- & 193 & 193 \\
\midrule
\multicolumn{4}{@{}l}{\emph{Humanities and Social Sciences}} \\
History & 695 & 252 & 947 \\
Geography & 473 & 134 & 607 \\
Philosophy & 102 & 254 & 356 \\
Sociology & 209 & 257 & 466 \\
Religious Studies & -- & 56 & 56 \\
Psychology & -- & 263 & 263 \\
Political Science & -- & 82 & 82 \\
International Relations & -- & 156 & 156 \\
Communication & -- & 300 & 300 \\
Library and Archival Science & -- & 482 & 482 \\
Social Work & -- & 95 & 95 \\
\bottomrule
\end{tabular}
\hfill
\begin{tabular}{@{}lrrr@{}}
\toprule
Subject & HS & UG & Total \\
\midrule
\multicolumn{4}{@{}l}{\emph{Education}} \\
Education & -- & 433 & 433 \\
\midrule
\multicolumn{4}{@{}l}{\emph{Law}} \\
Law & -- & 11,013 & 11,013 \\
\midrule
\multicolumn{4}{@{}l}{\emph{Business and Economics}} \\
Business Administration & -- & 1,586 & 1,586 \\
Public Administration & -- & 472 & 472 \\
Accounting & -- & 1,181 & 1,181 \\
Economics & -- & 585 & 585 \\
Finance & -- & 292 & 292 \\
\midrule
\multicolumn{4}{@{}l}{\emph{Health Sciences}} \\
Physical Education & 24 & 139 & 163 \\
Medicine & -- & 2,648 & 2,648 \\
Nursing & -- & 109 & 109 \\
Dentistry & -- & 98 & 98 \\
Pharmacy & -- & 84 & 84 \\
Nutrition & -- & 96 & 96 \\
Physical Therapy & -- & 119 & 119 \\
Speech and Language Therapy & -- & 102 & 102 \\
Occupational Therapy & -- & 45 & 45 \\
Biomedical Sciences & -- & 94 & 94 \\
Public Health & -- & 455 & 455 \\
\midrule
\multicolumn{4}{@{}l}{\emph{Agricultural and Veterinary Sciences}} \\
Agronomy & -- & 187 & 187 \\
Animal Science & -- & 136 & 136 \\
Veterinary Medicine & -- & 94 & 94 \\
\midrule
\multicolumn{4}{@{}l}{\emph{Services}} \\
Tourism and Hospitality & -- & 66 & 66 \\
Gastronomy & -- & 38 & 38 \\
Aesthetics and Cosmetology & -- & 36 & 36 \\
Occupational Health and Safety & -- & 111 & 111 \\
\bottomrule
\end{tabular}
\caption{The 72 (level, subject) pairs of taxonomy v1.0.0 and the questions
assigned to each. \emph{HS}: high school; \emph{UG}: undergraduate; ``--'':
the subject is not in that level's list. The smallest cells are high-school
Computer Science and Spanish (5 each), high-school Physical Education (24)
and undergraduate Aesthetics and Cosmetology (36) and Gastronomy (38).}
\label{tab:kc-subjects}
\end{table*}

% -----------------------------------------------------------------------------
\subsection{Prompt and decoding design}
\label{app:kc-prompt}

\paragraph{Motivation.} Two failure modes were anticipated. A model that
sees the exam, the edition or the course name may label a question by the
dominant subject of that source (every OAB question as Law, every medicine
course question as Medicine) instead of by its content; and a model that
reads thousands of exam texts may follow instructions embedded in them.
The prompt was written to prevent both, and the decoding to guarantee that
every record gets exactly one admissible label without a parsing step.

\paragraph{Prompt.} Each level has its own system message, built from a
fixed template: it states the task (classify multiple-choice questions by
their main discipline), asks for exactly one subject of the level's
taxonomy, and sets the criterion as \emph{the knowledge needed to answer the
command and to tell the choices apart}, adding that the narrative setting
and the language of the text do not determine the subject, that
interdisciplinary questions receive the central discipline, that the course,
job or exam must not be used as a shortcut, and that there is no ``Other'',
abstention or multi-subject answer. It then lists every allowed subject as
``name: definition. Includes: \ldots. Excludes: \ldots'', followed by the
thirteen boundary rules, and asks for the exact subject name in English, with
no explanation, JSON or macro-area. The user message is a JSON object with
the question and its choices indexed from zero, preceded by an instruction
to treat that content as data, even if it contains instructions. The
answer key, exam, edition, course, job title, URL and any earlier label are
not part of either message. Because the system message enumerates the
subjects, undergraduate prompts are about twice as long as high-school
prompts (4,137 versus 1,842 tokens for the shortest questions;
Table~\ref{tab:kc-cost}).

\paragraph{Decoding.} The model is Qwen3.5-122B-A10B-FP8, a 122B-parameter
mixture-of-experts model with 10B active parameters, in its official FP8
weights at revision \texttt{a099dee7}\ldots, served offline with
vLLM~0.19.1 (language model only, no expert parallelism, chat template
applied with the reasoning mode disabled). Generation uses temperature~0,
seed~42 and at most 32 output tokens, with a \emph{choice} structured-output
constraint over $T_\ell$, so the generated text is always one of the allowed
subject strings and no parsing or normalization is needed. The pipeline
then computes $\mu_\ell(s(x))$ and rejects the record if the subject is not
allowed for its level or if generation did not stop at the end of the
string. There is no retry and no default label: any such error aborts the
run. Each batch of 128 records is appended to a checkpoint with the SHA-256
of the full input row and the prompt and output token counts, and a run
resumes only when the input hash, model revision, generation parameters,
library versions, prompt, taxonomy and code hashes all match the stored
identity. Prompts whose tokens plus the 32-token answer budget would exceed
the context limit are rejected rather than truncated; the limit was set to
16,000 tokens and the longest prompt had 8,394.

\paragraph{What the mechanism guarantees.} All 37,887 records received a
label in one pass, every label is a valid pair $(\ell,s)$, every stored
macro-area equals $\mu_\ell(s)$, and all 72 pairs were used at least once.
Output lengths are two to seven tokens, essentially the tokenization of the
subject's name (only English varies, between two and four tokens, depending
on the preceding context). None of this says anything about whether the
subject is the right one; the constraint makes the label admissible, not
correct.

% -----------------------------------------------------------------------------
\subsection{Resulting distribution}
\label{app:kc-distribution}

\paragraph{By level and macro-area.} The two levels have different
profiles (Table~\ref{tab:kc-overview}). High-school questions are 56.8\%
STEM, 26.1\% Languages and Arts and 16.9\% Humanities and Social Sciences,
with 24 Physical Education questions as the only Health Sciences label.
Undergraduate questions are 37.8\% Law, a single subject, followed by STEM
(17.7\%), Business and Economics (14.1\%), Health Sciences (13.7\%),
Humanities and Social Sciences (8.0\%) and Languages and Arts (4.9\%);
Education, Agricultural and Veterinary Sciences and Services share the
remaining 3.8\%.

\paragraph{By subject.} Law is the largest subject with 11,013 questions,
followed by Medicine (2,648), Computer Science (2,171), Mathematics (2,153),
Business Administration (1,586), Portuguese (1,507), Logic (1,323) and
Accounting (1,181). The distribution has a long tail
(Figures~\ref{fig:kc-subjects-hs} and~\ref{fig:kc-subjects-ug}): 17 of the
72 pairs have fewer than 100
questions, and high-school Computer Science and Spanish have five each.
Within STEM, the two levels invert: Mathematics, Logic and Physics dominate
at high school, Computer Science and Engineering at undergraduate level.

\paragraph{By exam.} Table~\ref{tab:kc-exams} and Figure~\ref{fig:kc-exams}
show the labels of each exam. Seven exams have a declared scope, and their
labels fall almost entirely within it: 98.3\% of OAB and 98.1\% of ENAM are
Law; OBI is 91.3\% Logic and 8.4\% Mathematics, with three Computer Science
items; CFCES is 89.9\% Accounting and 9.8\% Law; REVALIDA and the USP/Unicamp
residency exams are 88.3\% and 87.1\% Medicine, with Public Health as the
second subject of REVALIDA (9.1\%); and POSCOMP is 68.6\% Computer Science
and 24.1\% Mathematics. The remaining exams are broad by design. ENADE uses
all 55 undergraduate subjects and all nine macro-areas, with Business
Administration first at only 12.3\%; the civil-service exams use 38 (BNDES),
29 (CNU) and 22 (BACEN) subjects; and the vestibular and military academy
exams spread over the school curriculum, with no subject above 20\% in ENEM,
FUVEST, COMVEST and BLUEX and a split between languages and STEM in AFA and
IME, whose first subject is English (24.4\% and 32.4\%).

\begin{table*}[t]
\centering
\footnotesize
\setlength{\tabcolsep}{4pt}
\begin{tabular}{@{}lrrrlrlr@{}}
\toprule
Exam & $n$ & Areas & Subjects & Most frequent subject & \% & Second subject & \% \\
\midrule
\multicolumn{8}{@{}l}{\emph{High school (8,746)}} \\
ENEM & 2,688 & 4 & 17 & Mathematics & 19.2 & Portuguese & 12.8 \\
FUVEST & 1,656 & 3 & 13 & History & 13.7 & Biology & 13.2 \\
OBI & 1,326 & 1 & 3 & Logic & 91.3 & Mathematics & 8.4 \\
AFA & 1,084 & 3 & 8 & English & 24.4 & Portuguese & 23.9 \\
IME & 879 & 3 & 9 & English & 32.4 & Mathematics & 26.5 \\
BLUEX & 661 & 3 & 13 & History & 18.2 & Mathematics & 16.2 \\
COMVEST & 452 & 3 & 13 & Mathematics & 18.1 & Physics & 13.9 \\
\midrule
\multicolumn{8}{@{}l}{\emph{Undergraduate (29,141)}} \\
OAB & 9,508 & 4 & 11 & Law & 98.3 & International Relations & 0.6 \\
ENADE & 9,237 & 9 & 55 & Business Administration & 12.3 & Computer Science & 5.6 \\
BNDES & 4,009 & 9 & 38 & Computer Science & 14.8 & Law & 11.8 \\
Res.\ USP/Unicamp & 1,491 & 5 & 19 & Medicine & 87.1 & Psychology & 3.8 \\
POSCOMP & 1,300 & 2 & 8 & Computer Science & 68.6 & Mathematics & 24.1 \\
REVALIDA & 1,195 & 5 & 10 & Medicine & 88.3 & Public Health & 9.1 \\
BACEN & 1,016 & 6 & 22 & Law & 34.7 & Computer Science & 13.6 \\
CFCES & 694 & 3 & 4 & Accounting & 89.9 & Law & 9.8 \\
CNU & 377 & 9 & 29 & Law & 17.5 & Public Administration & 17.2 \\
ENAM & 314 & 2 & 4 & Law & 98.1 & Sociology & 1.0 \\
\bottomrule
\end{tabular}
\caption{Labels per exam, sorted by size within each level. \emph{Areas} and
\emph{Subjects}: number of distinct macro-areas and subjects assigned to the
exam's questions; \emph{\%}: share of the exam's questions with the most
frequent and the second most frequent subject.}
\label{tab:kc-exams}
\end{table*}

\begin{figure*}[t]
\centering
\includegraphics[width=\textwidth]{\KnowledgeFigDir/macro_area_by_exam.pdf}
\caption{Share of each exam's questions assigned to every macro-area (\%).
Exams are grouped by academic level and sorted by size; empty cells are
macro-areas without any question of that exam, and ``$<$1'' marks shares
below one percent.}
\label{fig:kc-exams}
\end{figure*}

\begin{figure}[t]
\centering
\includegraphics[width=\columnwidth]{\KnowledgeFigDir/subjects_high_school.pdf}
\caption{Questions per subject at high-school level, grouped by macro-area
in the order of Table~\ref{tab:kc-subjects}.}
\label{fig:kc-subjects-hs}
\end{figure}

\begin{figure}[p]
\centering
\includegraphics[width=\columnwidth]{\KnowledgeFigDir/subjects_undergraduate.pdf}
\caption{Questions per subject at undergraduate level, grouped by
macro-area in the order of Table~\ref{tab:kc-subjects}. The scale is
linear, so Law (11,013) sets the axis and most subjects appear as short
bars; the count is printed next to each bar.}
\label{fig:kc-subjects-ug}
\end{figure}

% -----------------------------------------------------------------------------
\subsection{Consistency with withheld metadata}
\label{app:kc-audit}

\paragraph{Motivation.} No labelled reference exists for these questions,
so the audit cannot measure how often the assigned subject is right. It can
check whether the labels agree with information that was withheld from the
model and that constrains the plausible subjects of a question: the
disciplinary scope of an exam, the course of an ENADE booklet, or the
medical specialty of a residency programme. Agreement with such metadata is
weak supervision. It shows that the labels follow the content of exams and
courses rather than drifting at random, and it locates the groups where the
labels most often leave the expected set; it does not validate any
individual label.

\paragraph{Design.} Two kinds of groups are audited; the expected sets were
fixed by rule before the labels were inspected and were not adjusted
afterwards. (a)~Seven exams have a declared scope: Law for OAB and ENAM;
Accounting or Law for the accountants' exam (CFCES), whose questions on
corporate and procedural law are explicitly legal; Logic, Mathematics or
Computer Science for OBI; those three plus Statistics for POSCOMP; and
Medicine or Public Health for REVALIDA and the USP/Unicamp residency exams.
(b)~For ENADE, the edition string carries the course
(\texttt{ENADE 2013 - agronomia}; 471 editions, 151 course names). A list of
50 course families, each a pattern over course names, maps every course to
its home subject or subjects, adding Public Health for the health professions
and Education for teaching degrees (\emph{licenciaturas}, the pedagogy degree
and the former \emph{normal superior}). The patterns cover 148 courses and
8,875 questions (96.1\% of ENADE); three technology courses (hospital
management, agribusiness and radiology) have no single home subject and are
left unaudited. ENADE booklets have a general-education component shared by
all courses of a year followed by a course-specific component, so we split the
questions at the printed number: questions 1--10 are taken as general and
questions 11 and above as course-specific. The general component should not
favour the course's subjects and serves as a contrast. Mixed-content exams
(BNDES, BACEN, CNU and the vestibular and military exams) have no expected
set and are not audited.

\paragraph{Results for exams.} Table~\ref{tab:kc-audit}(a) lists the
seven exams. Between 89.5\% and 100\% of their labels fall inside the
expected set: all 1,326 OBI questions are Logic, Mathematics or Computer
Science, and CFCES, POSCOMP, OAB and ENAM are above 98\%. The labels outside
the set are neighbours of the scope rather than unrelated fields. In OAB, 55
questions on public international law (international organizations, dispute
settlement, diplomatic immunities) are International Relations, and 42 are
Philosophy; in the residency exams, 57 questions are Psychology, 49 of them
from the psychotherapy and psychiatry programmes, and 27 are Law; in
REVALIDA, 11 clinical scenarios that turn on legal or ethical duties are Law.
These are exactly the boundaries that rules 4, 7 and 12 of
\S\ref{app:kc-taxonomy} try to settle, and the audit shows where the model
drew them.

\paragraph{Results for ENADE.} Over the course-specific component of the
148 mapped courses, 68.7\% of the labels (5,732 of 8,347) fall inside the
course's expected set, and the most frequent label outside it is Business
Administration (281). The share varies widely between families
(Table~\ref{tab:kc-audit}(b), Figure~\ref{fig:kc-enade}). Among families with
at least 140 questions it is highest for Medicine (94.3\%), Law (88.9\%),
Philosophy (87.6\%), Mathematics (84.7\%), Pedagogy (83.0\%) and Physical
therapy (82.6\%), and lowest for Tourism (41.4\%), Engineering (44.9\%),
Agronomy and forestry (50.4\%), Social work (51.5\%) and Environmental
engineering (57.0\%); over all 50 families it ranges from 34.2\%
(Biomedicine) to 95.9\% (Archival and library science). The subject most
often found outside the expected set is in every case an adjacent discipline:
Business Administration in the engineering, tourism, computing, design and
communication courses; Biology in biomedicine (35 of 111) and agronomy (37
of 230); Sociology in social work (45 of 167); Accounting in financial
management (41 of 118); Food Science and Technology in gastronomy (33 of 87)
and nutrition (30 of 152); Medicine in pharmacy, dentistry, nursing and
physical therapy; and Education in the teaching degrees, where 44.5\% of the
course-specific questions are labelled Education and the rest with the
discipline taught, as rule~3 prescribes. In the general component, 58.5\% of
the labels (309 of 528) fall inside the course's set, 10~percentage points
less than in the specific component, and its most frequent subjects are
Business Administration (66), Sociology (27), Public Administration (24) and
Public Health (24), which fits the civic and organizational themes of that
component. Two caveats apply to this contrast. The general component is
small, 545 questions with a median of two per edition, because exact and
semantic deduplication keep a single copy of a question shared by all
courses of a year, and the surviving copy is attributed to one course by the
deduplication order, not by content. And the split at question~10 is a
heuristic: in some editions the course-specific component starts at
question~1, which is visible in families whose general component agrees with
the course as strongly as the specific one.

\begin{table*}[t]
\centering
\footnotesize
\setlength{\tabcolsep}{4pt}
\begin{tabular}{@{}lrp{0.34\textwidth}rrl@{}}
\toprule
Group & $n$ & Expected subjects & Within (\%) & 95\% CI & Most frequent other ($n$) \\
\midrule
\multicolumn{6}{@{}l}{\emph{(a) Exams with a declared scope}} \\
OBI & 1,326 & Logic, Mathematics, Computer Science & 100.0 & 99.7--100.0 & -- \\
CFCES & 694 & Accounting, Law & 99.7 & 99.0--99.9 & Finance (1) \\
POSCOMP & 1,300 & Computer Science, Mathematics, Statistics, Logic & 99.5 & 98.9--99.7 & Engineering (3) \\
OAB & 9,508 & Law & 98.3 & 98.0--98.6 & International Relations (55) \\
ENAM & 314 & Law & 98.1 & 95.9--99.1 & Sociology (3) \\
REVALIDA & 1,195 & Medicine, Public Health & 97.4 & 96.3--98.2 & Law (11) \\
Res.\ USP/Unicamp & 1,491 & Medicine, Public Health & 89.5 & 87.9--91.0 & Psychology (57) \\
\midrule
\multicolumn{6}{@{}l}{\emph{(b) ENADE, course-specific component (questions 11 and above), families with at least 140 questions}} \\
Business and management & 958 & Business Administration & 75.7 & 72.9--78.3 & Law (47) \\
Engineering & 839 & Engineering & 44.9 & 41.6--48.3 & Business Administration (71) \\
Computing & 566 & Computer Science$^\dagger$ & 80.2 & 76.7--83.3 & Business Administration (37) \\
Agronomy and forestry & 230 & Agronomy & 50.4 & 44.0--56.8 & Biology (37) \\
Communication & 207 & Communication & 63.3 & 56.5--69.6 & Business Administration (29) \\
Physical education & 204 & Physical Education, Public Health$^\dagger$ & 74.5 & 68.1--80.0 & Education (32) \\
Design & 203 & Design & 72.4 & 65.9--78.1 & Business Administration (26) \\
Environmental engineering & 193 & Environmental Sciences, Engineering & 57.0 & 49.9--63.8 & Law (18) \\
Social work & 167 & Social Work & 51.5 & 44.0--59.0 & Sociology (45) \\
Animal science & 161 & Animal Science & 64.6 & 56.9--71.6 & Agronomy (21) \\
Philosophy & 161 & Philosophy$^\dagger$ & 87.6 & 81.6--91.8 & Logic (11) \\
Nursing & 155 & Nursing, Public Health & 81.9 & 75.1--87.2 & Medicine (15) \\
Nutrition & 152 & Nutrition, Public Health & 62.5 & 54.6--69.8 & Food Science and Technology (30) \\
Food technology & 151 & Food Science and Technology & 73.5 & 66.0--79.9 & Engineering (10) \\
Veterinary medicine & 148 & Veterinary Medicine & 58.1 & 50.1--65.8 & Animal Science (20) \\
Tourism & 145 & Tourism and Hospitality & 41.4 & 33.7--49.5 & Business Administration (34) \\
Physical therapy & 144 & Physical Therapy, Public Health & 82.6 & 75.6--88.0 & Medicine (16) \\
Mathematics & 144 & Mathematics$^\dagger$ & 84.7 & 78.0--89.7 & Education (12) \\
Law & 144 & Law & 88.9 & 82.7--93.0 & Philosophy (5) \\
Dentistry & 144 & Dentistry, Public Health & 78.5 & 71.1--84.4 & Medicine (21) \\
Accounting & 142 & Accounting & 77.5 & 69.9--83.6 & Business Administration (9) \\
Arts & 142 & Arts$^\dagger$ & 71.8 & 63.9--78.6 & Education (25) \\
Psychology & 141 & Psychology & 81.6 & 74.4--87.1 & Sociology (11) \\
Pedagogy & 141 & Education & 83.0 & 75.9--88.3 & Philosophy (6) \\
Economics & 140 & Economics & 72.9 & 65.0--79.5 & Statistics (19) \\
Medicine & 140 & Medicine, Public Health & 94.3 & 89.1--97.1 & Psychology (2) \\
24 smaller families & 2,285 & (per family) & 67.8 & 65.9--69.7 & -- \\
All 50 families & 8,347 & (per family) & 68.7 & 67.7--69.7 & Business Administration (281) \\
Not audited (3 courses) & 362 & -- & -- & -- & -- \\
\midrule
\multicolumn{6}{@{}l}{\emph{(c) ENADE, general-education component (questions 1--10)}} \\
All mapped courses & 528 & (per family) & 58.5 & 54.3--62.6 & Sociology (25) \\
\bottomrule
\end{tabular}
\caption{Agreement between the assigned subjects and the subjects expected
from withheld metadata. \emph{Within}: percentage of the group's questions
whose subject belongs to the expected set, with its 95\% Wilson interval;
\emph{Most frequent other}: the most frequent subject outside the expected
set, with its count. $^\dagger$Education is also expected for the
\emph{licenciatura} courses of the family. Families are sorted by size; the
24 families with fewer than 140 questions are aggregated. The expected sets
were fixed by rule before inspecting the labels; a label outside the set is
not an error, and a label inside it is not confirmed.}
\label{tab:kc-audit}
\end{table*}

\begin{figure}[t]
\centering
\includegraphics[width=\columnwidth]{\KnowledgeFigDir/enade_course_consistency.pdf}
\caption{Share of ENADE labels inside the course's expected set, per course
family with at least 30 course-specific questions, for the course-specific
component (filled, with 95\% Wilson interval) and the general-education
component (open, shown when it has at least 10 questions). Labels at the
right give the two group sizes.}
\label{fig:kc-enade}
\end{figure}

\paragraph{Reading.} The shares above measure consistency, not accuracy.
The expected sets describe a course or an exam, not a question: a question in
a law course may require Philosophy or Economics, and a question in an
engineering course may be about management, so a label outside the set is
not an error, and a label inside it is not confirmed. The audit does show
that the labels track content where content and metadata are strongly
coupled, that they do not collapse to the exam's dominant subject (the model
never saw the exam), and that the groups with the lowest agreement are those
whose home subject has close neighbours in the taxonomy (biomedicine and
biology, finance and accounting, gastronomy and food technology, tourism and
business administration). Attributing the disagreements to the model, to the
taxonomy or to the questions themselves would require reading the items.

% -----------------------------------------------------------------------------
\subsection{Illustrative boundary cases}
\label{app:kc-boundary}

Table~\ref{tab:kc-boundary} lists questions sampled deterministically
(seed~42) from the groups where the boundary rules apply: the two sides of
each boundary within the exams where it matters most, and the non-Law labels
of OAB. They illustrate what each label looks like in practice; no judgement
of correctness was made, and the observations describe the item, not the
decision.

\begin{table*}[t]
\centering
\footnotesize
\setlength{\tabcolsep}{3pt}
\begin{tabular}{@{}llp{0.40\textwidth}lp{0.24\textwidth}@{}}
\toprule
Boundary & Exam, edition, item & Opening words & Subject & Observation \\
\midrule
Portuguese / Literature & ENEM 2021, q.~44 & \emph{Naquele tempo, Itaguaí, que, como as demais vilas, arraiais e povoações da colônia, não dispunha de imprensa\ldots} & Portuguese & Excerpt from a Brazilian novel, read as a text. \\
 & ENEM 2021, q.~39 & \emph{Falso moralista. Você condena o que a moçada anda fazendo e não aceita o teatro de revista\ldots} & Portuguese & Song lyrics in a colloquial register. \\
 & ENEM 2020, q.~4 & \emph{Oye, Pito, ésta es: la vida bruta de un boy. Mis tierras eran Nuevo México, Colorado, California\ldots} & Literature & Poem in Spanish from the foreign-language block; labelled Literature, not Spanish. \\
\midrule
Logic / Mathematics / Computer Science & OBI 2008, phase 2, level 1, q.~16 & \emph{Viagem a Marte. Três engenheiros \ldots devem enviar exatamente cinco astronautas para Marte dos oito que\ldots} & Logic & Selection under constraints given in the statement. \\
 & OBI 2015, phase 1, level 2, q.~3 & \emph{O Rei da Nlogônia decidiu organizar um torneio de tênis com os dez melhores jogadores do reino. Inicialmente cada jogador ganha uma moeda de ouro\ldots} & Mathematics & Tournament puzzle with coins exchanged between players. \\
 & OBI 2006, phase 2, level 2, q.~31 & \emph{Representação Pós-Fixa. Desde pequenos aprendemos a escrever expressões aritméticas\ldots} & Computer Science & Postfix notation of arithmetic expressions. \\
\midrule
Computer Science / Mathematics / Logic / Statistics & POSCOMP 2012, q.~47 & \emph{O fenômeno de thrashing de um sistema é caracterizado por:} & Computer Science & Operating-systems term. \\
 & POSCOMP 2012, q.~16 & \emph{Com relação à proposição $P$: ``Seja $a\in\mathbb{N}$. Se $a^2$ é ímpar então $a$ é ímpar''\ldots} & Mathematics & Proof techniques for a number-theoretic statement; labelled Mathematics, not Logic. \\
 & POSCOMP 2007, q.~11 & \emph{Considere a seguinte linguagem de primeira ordem: constantes $a, b$; variáveis $x, y$; predicados unários $P$\ldots} & Logic & First-order logic. \\
 & POSCOMP 2023, q.~19 & \emph{Com base na Tabela 1, calcule a média da seguinte amostra de número de filhos\ldots} & Statistics & Sample mean from a frequency table. \\
\midrule
Accounting / Law & CFCES, 16th QTG, q.~10 & \emph{Conforme a NBC TA 570 -- Continuidade Operacional, os objetivos do auditor independente são, EXCETO:} & Accounting & Auditing standard cited; labelled Accounting. \\
 & CFCES, 10th QTG, q.~5 & \emph{De acordo com a Lei n.º 6.404/76, a companhia poderá negociar com as próprias ações, EXCETO:} & Law & Corporate-law provision asked directly. \\
 & CFCES, 5th forensic accounting, q.~16 & \emph{Com base no que dispõe o Código de Processo Civil, analise os itens abaixo\ldots No laudo, o perito deve apresentar\ldots} & Law & Civil-procedure rules on expert reports. \\
\midrule
Medicine / Public Health & REVALIDA 2021, q.~15 & \emph{Um paciente, 2 anos, sexo masculino, chega ao pronto atendimento \ldots com relato de ter iniciado há 5 dias coriza serosa e tosse seca\ldots} & Medicine & Clinical case of one patient. \\
 & REVALIDA 2016, q.~71 & \emph{Um município de 15 mil habitantes deseja cobrir 100\% do seu território com equipes de Saúde da Família\ldots} & Public Health & Health-system planning for a population. \\
\midrule
Education / taught discipline & ENADE 2014, computing (lic.), q.~26 & \emph{A elaboração e a disponibilização de material didático interativo é importante para a motivação do estudante\ldots} & Education & Design of teaching material. \\
 & ENADE 2014, philosophy (lic.), q.~26 & \emph{Ora, é através da aprendizagem de alguns saberes práticos (savoir-faire) envolvidos na inculcação massiva da ideologia da classe dominante\ldots} & Sociology & Theoretical passage on ideology and schooling; labelled Sociology, neither Education nor Philosophy. \\
 & ENADE 2014, history (lic.), q.~23 & \emph{A escassez de testemunhos sobre o comportamento e atitudes das classes subalternas do passado\ldots} & History & Historiography; labelled with the discipline taught. \\
\midrule
Non-Law labels in OAB & OAB-MG 2007.3, q.~97 & \emph{São funções e competências do Conselho de Segurança da Organização das Nações Unidas (ONU) na forma prevista pela Carta das Nações Unidas de 1945, exceto:} & International Relations & Competences of a UN body under the UN Charter; labelled International Relations rather than Law. \\
\bottomrule
\end{tabular}
\caption{Items sampled with seed~42 from the boundary groups (two per
subject and group; the script prints the full sample, of which 19 items are
listed). \emph{Opening words}: beginning of the question, lightly shortened;
\emph{Subject}: the assigned label; \emph{Observation}: a description of the
item, not an assessment of the label. ``lic.'': \emph{licenciatura}
(teaching degree); ``QTG'': general technical qualification exam.}
\label{tab:kc-boundary}
\end{table*}

% -----------------------------------------------------------------------------
\subsection{Cost and reproducibility}
\label{app:kc-cost}

Table~\ref{tab:kc-cost} summarizes the run. The 37,887 questions were
classified on one NVIDIA B200 in 84.9~minutes end to end (5,092.8~s,
measured from the start of the worker process to the end of the last batch,
and therefore including 35.9~s to load the 115.4~GiB of FP8 weights and
155.9~s to initialize the engine, of which 122.9~s are the warm-up run), at
7.44 questions per second; between the first and the last batch the rate was
7.84 questions per second. The engine held a KV cache of 454,832 tokens
(41.7~GiB), enough for 79 concurrent requests at the 16,000-token context
limit, and ran up to 32 sequences at a time in batches of 128 records (296
rounds). The run processed 145.6~million prompt tokens, 2,110.5 per
high-school question and 4,362.8 per undergraduate question on average; the
fixed part of the prompt dominates, since the shortest prompts have 1,842 and
4,137 tokens and the 99th percentiles are 3,134 and 4,974. Output totalled
101,739 tokens, 2.7 per question. Prefix caching was disabled, so the
prompt-token count is the full cost of reprocessing the shared system message
for every question and an upper bound for a re-run with caching. The run
manifest records the model identifier and resolved revision, the resolved
input hash, the generation parameters, the library versions (Python~3.12.3,
vLLM~0.19.1, PyTorch~2.10.0, Transformers~4.57.6, Datasets~5.0.1), the SHA-256
of the prompt template (\texttt{8ffbe0c2}\ldots), of the taxonomy
(\texttt{3b2cdd5e}\ldots) and of the implementation (\texttt{cd9872d9}\ldots),
and the SHA-256 of the checkpoint (\texttt{fbab3655}\ldots). Temperature~0
with a fixed seed makes the run deterministic on the same hardware and
software stack, but identical labels are not guaranteed across GPUs, kernels
or parallelism settings.

\begin{table}[t]
\centering
\footnotesize
\setlength{\tabcolsep}{4pt}
\begin{tabular}{@{}lr@{}}
\toprule
Hardware and engine & \\
\midrule
GPU & 1$\times$ NVIDIA B200, 191.5~GB \\
Weights (FP8) & 115.37~GiB, loaded in 35.9~s \\
Engine initialization & 155.9~s (warm-up 122.9~s) \\
KV cache & 454,832 tokens, 41.7~GiB \\
Context limit / concurrent sequences & 16,000 tokens / 32 \\
Batch size / rounds & 128 / 296 \\
Data / tensor parallelism & 1 / 1 \\
Prefix caching / chunked prefill & off / on \\
End-to-end inference & 5,092.8~s (84.9~min) \\
Throughput, end to end / between batches & 7.44 / 7.84 q/s \\
\bottomrule
\end{tabular}

\medskip
\begin{tabular}{@{}lrrr@{}}
\toprule
Tokens & High school & Undergraduate & All \\
\midrule
Questions & 8,746 & 29,141 & 37,887 \\
Prompt, total & 18,458,112 & 127,137,386 & 145,595,498 \\
Prompt, mean & 2,110.5 & 4,362.8 & 3,842.9 \\
Prompt, median & 2,052 & 4,327 & 4,280 \\
Prompt, minimum & 1,842 & 4,137 & 1,842 \\
Prompt, 95th percentile & 2,539 & 4,649 & 4,604 \\
Prompt, 99th percentile & 3,134 & 4,974 & 4,913 \\
Prompt, maximum & 3,759 & 8,394 & 8,394 \\
Output, total & 23,825 & 77,914 & 101,739 \\
\bottomrule
\end{tabular}
\caption{Cost of the production run. Top: hardware, engine settings and
wall-clock times from the run log and manifest. Bottom: prompt and output
tokens recorded per record in the run checkpoint, by academic level.}
\label{tab:kc-cost}
\end{table}

% -----------------------------------------------------------------------------
\subsection{Limitations}
\label{app:kc-limitations}

\emph{A single pass.} Model, prompt, taxonomy and decoding were fixed before
the run and never varied, so nothing in this appendix shows that another
model, a shorter prompt, a different taxonomy or reasoning-enabled decoding
would have produced different or worse labels.
\emph{No accuracy.} Without a labelled reference, the audit measures only
agreement with group-level metadata. It cannot say how many individual
labels are right, and agreement with a course or exam is compatible with
systematic confusions between neighbouring subjects.
\emph{Expected sets are a judgement.} The audit's expected sets were fixed by
a stated rule before the labels were inspected, but the rule itself (home
subjects, Public Health for health professions, Education for teaching
degrees) is the authors' choice, and the general-education component of ENADE
is identified by a question-number heuristic that fails in some editions.
\emph{One label per question.} Interdisciplinary questions receive a single
subject and no question can be marked as unclassifiable; the constrained
decoding guarantees an admissible label even where the taxonomy fits
poorly.
\emph{Greedy decoding under a grammar.} The label is the greedy token path
among the allowed strings, which need not be the subject of highest total
probability when two subjects share a prefix (Physical Education and
Physical Therapy; Public Administration and Public Health); the effect was
not measured.
\emph{A project-specific taxonomy.} The lists were designed for these 17
exams and do not follow an official classification of fields of study;
results per subject are not directly comparable to other benchmarks.
\emph{Imbalance.} Law alone holds 29\% of the questions, 17 of the 72
(level, subject) pairs have fewer than 100, and two have five; per-area
statistics and stratified sampling must account for this.
\emph{Reproducibility.} Deterministic decoding does not guarantee identical
labels across GPUs, kernels or parallelism settings, and the reported cost is
an upper bound, since prefix caching was disabled.
